\section{Related Work}
\label{sec::related}

\subsection{Simulation in Macroeconomics}\label{sec::relatedwork1}

ABM emerged as a more promising solution for macroeconomic research compared with empirical statistical models~\cite{hendry1982formulation,phelps1967phillips,kydland1982time} and DSGE models~\cite{christiano2005nominal} in recent decades.
Diverse agents interact with each other based on rules or computational models, avoiding the assumption of a predetermined economic equilibrium. These models allow for a wide range of nonlinear behaviors, enabling policymakers to simulate different policy scenarios and qualitatively assess their impacts on the economy.
The agent modeling with predetermined rules~\cite{tesfatsion2006handbook,brock1998heterogeneous} or neural networks~\cite{trott2021building,zheng2022ai,mi2023taxai} still suffers from the assumption of oversimplified agent behavior or the dependence on large-scale training data, thus having limitations in capturing the full complexity of economic systems.

In this work, we introduce EconAgent with reasoning and planning abilities for simulating macroeconomic activities.



\subsection{LLM-empowered Agents}\label{sec::relatedwork1}
Recently, LLMs trained with large-scale corpus have shown human-level abilities and provide the basis for constructing agents for simulation~\cite{wang2023survey,xi2023rise}.
The LLM agents primarily have three advantages for simulation, including autonomous but adaptive reactions~\cite{autogpt, babyagi}, human-like intelligence for planning~\cite{wang2023survey,xi2023rise}, and interaction and communication with other agents or human~\cite{park2023generative,gilbert2005simulation,park2023generative}. With these advanced abilities, LLM agents have been widely used in many areas, including social science~\cite{park2022social,park2023generative,kovavc2023socialai,gao2023s,jinxin2023cgmi}, natural science~\cite{boiko2023emergent,bran2023chemcrow}, \textit{etc}. Moreover, some works also consider LLM agents in economic research, which can be categorized into three levels~\cite{gao2023large}, including rationality or bias in individual behavior~\cite{horton2023large,chen2023emergence}, planning and cooperation in interactive behavior~\cite{guo2023gpt,akata2023playing}, and system-level market simulation~\cite{zhao2023competeai,anonymous2024rethinking,chen2023put}. 

In short, existing works only consider individual one-step or few-step behavior for a few agents, without experimenting on multi-step behaviors within a multi-agent simulation environment, which is the focus of our work.

